\chapter{Curvature Lines, Ridges and Valleys}
\label{ch:vii42}

The smooth theory of VII/16--VII/17 associates to every nonumbilic point two
principal directions and two principal curvatures.  Those directions form line
fields whose integral curves are curvature lines.  In applications, derivatives
of the principal curvatures along these lines identify ridge and valley
features.  This final chapter joins the smooth definitions to the estimation,
scale selection, and tracing problems that arise on triangle meshes.

\section{Principal directions revisited}

Let \(S:T_pM\to T_pM\) be the shape operator of an oriented smooth surface.  At
a nonumbilic point it has distinct eigenvalues \(k_1,k_2\) and orthogonal unit
eigenvectors \(e_1,e_2\):
\[
Se_i=k_i e_i,
\qquad
\langle e_1,e_2\rangle=0.
\]
We locally order the eigenvalues as \(k_1>k_2\) when convenient.

\begin{warning}[Directions are lines, not arrows]\label{warn:vii42-linefield}
The eigenvectors \(e_i\) and \(-e_i\) represent the same principal direction.
A principal-direction field is therefore naturally a line field.  Numerical
algorithms that treat it as an ordinary oriented vector field can introduce
spurious sign flips.
\end{warning}

\section{Euler's normal-curvature formula}

If a unit tangent vector makes angle \(\theta\) with \(e_1\), Euler's formula is
\[
k_n(\theta)=k_1\cos^2\theta+k_2\sin^2\theta.
\]
Thus the principal directions are exactly the directions in which normal
curvature is extremal.  This gives a geometric interpretation of the eigenvectors
that does not depend on a particular parameterization.

\section{Curvature lines}

\begin{definition}[Line of curvature]\label{def:vii42-curvature-line}
A regular surface curve \(\gamma\) is a line of curvature if its tangent is a
principal direction at every nonumbilic point of the curve:
\[
\gamma'(s)\parallel e_i(\gamma(s))
\]
for one principal line field \(i\in\{1,2\}\).
\end{definition}

Away from umbilics and points where the chosen field loses regularity, ordinary
differential-equation theory gives local integral curves.  Since \(e_1\perp
e_2\), the two curvature-line families cross orthogonally wherever the
principal curvatures are distinct.

\begin{proposition}[Normal derivative along a curvature line]\label{prop:vii42-normal}
If \(\gamma\) is a unit-speed curvature line with tangent \(e_i\), then, under
the convention \(S=-dN\),
\[
\frac{d}{ds}N(\gamma(s))=-k_i e_i.
\]
\end{proposition}

\section{Umbilics and singular line fields}

At an umbilic,
\[
k_1=k_2,
\qquad
S=kI,
\]
so every tangent direction is principal and no distinguished principal line
field exists.  Umbilics are therefore genuine singularities for curvature-line
tracing, not merely numerical nuisances.

On a round sphere every point is umbilic.  On a circular cylinder there are no
umbilics: one principal family follows the generators and the other follows the
circles.

\begin{remark}[Topology enters through singularities]\label{rem:vii42-topology}
On a closed surface, principal line fields cannot generally be made globally
nonsingular.  The arrangement and indices of umbilic singularities are tied to
global topology, foreshadowing the homotopy and homology viewpoint of Volume
VIII.
\end{remark}

\section{Directional derivatives of principal curvature}

For a principal curvature \(k_i\), define
\[
D_{e_i}k_i=\langle\nabla k_i,e_i\rangle.
\]
A zero of this derivative means that \(k_i\) is stationary to first order when
moving in its own principal direction.  This is the basic differential signal
used by many ridge and valley definitions.

\section{A declared ridge--valley convention}

Terminology varies across geometry processing and computer vision, so we make a
specific local convention.

\begin{definition}[Principal ridge candidate]\label{def:vii42-ridge}
On an oriented patch with \(k_1>k_2\), a \(k_1\)-ridge candidate is a nonumbilic
point satisfying
\[
D_{e_1}k_1=0.
\]
It is a strict ridge point for this convention when additionally
\[
D_{e_1}^2k_1<0.
\]
\end{definition}

\begin{definition}[Principal valley candidate]\label{def:vii42-valley}
A \(k_2\)-valley candidate satisfies
\[
D_{e_2}k_2=0,
\]
and is strict when
\[
D_{e_2}^2k_2>0.
\]
\end{definition}

These definitions encode maxima of the larger principal curvature along its
principal direction and minima of the smaller principal curvature along its
principal direction.

\begin{warning}[Orientation and convention dependence]\label{warn:vii42-orientation}
Reversing the surface normal changes \(S\) to \(-S\) and flips principal
curvature signs.  It can also change how an ordering such as \(k_1>k_2\) is
named.  A ridge/valley implementation must therefore document its normal,
ordering, and sign conventions.  Some applications instead use absolute
curvature or other orientation-insensitive feature measures.
\end{warning}

\section{Ridges are not curvature lines}

A curvature line is an integral curve of a principal direction field.  A ridge
or valley is a locus where a directional derivative of principal curvature
vanishes and a second-derivative test selects an extremum.  A ridge curve may be
traced across the surface, but its tangent need not itself equal the principal
direction used in the defining derivative.

\begin{warning}[Common conceptual error]\label{warn:vii42-ridge-line}
Do not identify ``ridge curve'' with ``line of curvature.''  The first is an
implicit feature locus; the second is an integral curve of an eigenline field.
They coincide only under additional geometry.
\end{warning}

\section{Estimating curvature on triangle meshes}

A triangle mesh has no classical pointwise second fundamental form at a vertex.
One therefore estimates smooth local geometry.  A common pipeline is:
\[
\text{neighborhood}
\to
\text{normal/tangent frame}
\to
\text{local surface fit}
\to
\text{shape operator}
\to
(k_i,e_i).
\]

For example, in a local tangent frame \((x,y,z)\), fit a Monge patch
\[
z(x,y)
\approx
\frac12\bigl(ax^2+2bxy+cy^2\bigr)
+
\text{lower/higher-order terms}.
\]
At small slope, the Hessian
\[
\begin{pmatrix}a&b\\b&c\end{pmatrix}
\]
approximates the shape operator in that frame.  More complete fits account for
first derivatives, metric effects, and weighting.

\section{Neighborhood and scale selection}

Curvature is a second-derivative quantity and is therefore sensitive to noise.
A one-ring fit may be too noisy, while a very large neighborhood can erase
narrow features.  Robust implementations expose a physical or intrinsic scale,
use weighted least squares, reject degenerate neighborhoods, and compare
features across multiple scales.

\begin{definition}[Umbilic tolerance]\label{def:vii42-umbilic-tol}
A practical tracer treats a vertex as numerically near-umbilic when
\[
|k_1-k_2|<\varepsilon_k
\]
for a scale-aware tolerance \(\varepsilon_k\).
\end{definition}

Near this set, principal directions are ill-conditioned: a small perturbation of
the estimated shape operator can rotate the eigenvectors substantially.

\section{Cohering a principal line field}

Since \(e\) and \(-e\) are equivalent, neighboring estimates should be compared
by their unoriented angle.  When a temporary orientation is required for
interpolation, choose the sign that maximizes the dot product with a previously
oriented neighbor.  This creates local coherence but must not be mistaken for a
global orientation across line-field singularities.

\section{Extracting ridge and valley loci}

Suppose a scalar estimate
\[
r_1=D_{e_1}k_1
\]
is available at mesh vertices.  A standard first step is to detect sign changes
of \(r_1\) along edges and linearly interpolate candidate zero crossings within
triangles.  Candidates are then filtered by the second-directional-derivative
criterion, curvature strength, confidence, and proximity to umbilics.

\begin{remark}[Zero crossing is only a candidate test]\label{rem:vii42-zero}
A sign change in a noisy first derivative can create a false feature.  Stable
feature extraction generally combines derivative tests with strength and
persistence thresholds and with geometric continuity across neighboring faces.
\end{remark}

\section{Tracing curvature lines on a mesh}

To trace a curvature line, start from a seed and integrate an estimated principal
line field in the tangent plane.  Within a triangle, interpolate a representation
that respects the sign ambiguity, advance to the next edge, transport the local
frame as needed, and stop at boundaries, singular regions, low-confidence zones,
or after closing a loop.

Naively interpolating signed eigenvectors can cause cancellation when two
vertices store opposite representatives of the same line.  Line-field-aware
interpolation or local sign alignment is essential.

\section{Canonical examples}

\begin{example}[Round sphere]\label{ex:vii42-sphere}
For a round sphere, \(k_1=k_2\) everywhere.  Every direction is principal, so a
principal-direction tracer has no unique field to follow.  The sphere is a
useful regression test for correct umbilic detection.
\end{example}

\begin{example}[Circular cylinder]\label{ex:vii42-cylinder}
For a cylinder of radius \(R\), the principal curvatures are \(1/R\) and \(0\)
up to orientation.  Their principal directions are circumferential and axial.
Both curvatures are constant, so the directional derivative ridge tests vanish
everywhere; this illustrates why derivative zero alone is insufficient for
selecting salient ridges.
\end{example}

\begin{example}[Noisy scan]\label{ex:vii42-noisy}
On a noisy triangulated scan, raw one-ring curvature derivatives can oscillate in
sign at nearly every edge.  Increasing the fit scale or smoothing the estimated
shape operator can reveal persistent feature curves that survive across a range
of scales.
\end{example}

\section{Validation and regression geometry}

Feature-line code should be tested on analytic surfaces with known principal
geometry before it is trusted on arbitrary meshes.  Useful tests include a
plane, sphere, cylinder, torus patches, and synthetic Monge surfaces.  Numerical
checks should distinguish errors in normals, curvature magnitudes, principal
directions, directional derivatives, and final line extraction.

For a direction estimate, compare lines rather than oriented vectors.  If
\(e\) is estimated and \(e_*\) exact, a suitable error is
\[
\theta_{\rm line}
=
\arccos\bigl(|\langle e,e_*\rangle|\bigr),
\qquad
0\le\theta_{\rm line}\le\frac\pi2.
\]

\section{Relationship to the discrete Laplacian}

VII/41 showed that \(M^{-1}KX\) supplies a mean-curvature-normal estimate.
Principal curvatures require more information: mean curvature alone gives only
\(k_1+k_2\).  Recovering directions and separating \(k_1\) from \(k_2\) needs a
shape-operator or equivalent second-order estimate.  The Laplacian is therefore
one useful ingredient, not a complete ridge detector.

\section{A production feature-line pipeline}

A robust pipeline typically performs the following conceptual steps:
\[
\begin{aligned}
&\text{mesh validation and scale choice}\\
&\to \text{normal and tangent-frame estimation}\\
&\to \text{local shape-operator fitting}\\
&\to \text{principal curvature/direction estimation}\\
&\to \text{line-field coherence and umbilic masking}\\
&\to \text{directional curvature derivatives}\\
&\to \text{zero-crossing candidates and extremum tests}\\
&\to \text{connectivity, pruning, and multiscale validation}.
\end{aligned}
\]
Each stage should expose confidence diagnostics rather than silently passing
ill-conditioned data downstream.


\section{Exercises}

\begin{exercise}\label{exr:vii42-01}
What are principal directions at a nonumbilic point?
\end{exercise}

\begin{hint}
Use the shape operator.
\end{hint}

\begin{solution}
They are the two orthogonal eigenlines of the self-adjoint shape operator \(S\), with eigenvalues equal to the principal curvatures.
\end{solution}

\begin{exercise}\label{exr:vii42-02}
Why is a principal direction naturally a line rather than an oriented vector?
\end{exercise}

\begin{hint}
Compare \(e\) and \(-e\).
\end{hint}

\begin{solution}
If \(Se=ke\), then \(S(-e)=k(-e)\); the two vectors represent the same one-dimensional eigenspace.
\end{solution}

\begin{exercise}\label{exr:vii42-03}
State Euler's formula for normal curvature.
\end{exercise}

\begin{hint}
Measure angle from \(e_1\).
\end{hint}

\begin{solution}
\(k_n(\theta)=k_1\cos^2\theta+k_2\sin^2\theta\).
\end{solution}

\begin{exercise}\label{exr:vii42-04}
Define a line of curvature.
\end{exercise}

\begin{hint}
Its tangent follows an eigendirection.
\end{hint}

\begin{solution}
It is a regular surface curve whose tangent is parallel to one principal direction at each nonumbilic point.
\end{solution}

\begin{exercise}\label{exr:vii42-05}
How do the two curvature-line families meet at a nonumbilic point?
\end{exercise}

\begin{hint}
The shape operator is self-adjoint.
\end{hint}

\begin{solution}
They meet orthogonally because eigenvectors for distinct eigenvalues of a self-adjoint operator are orthogonal.
\end{solution}

\begin{exercise}\label{exr:vii42-06}
What happens to principal directions at an umbilic?
\end{exercise}

\begin{hint}
The eigenvalues coincide.
\end{hint}

\begin{solution}
The shape operator is a scalar multiple of the identity, so every tangent direction is principal and no distinguished principal line field exists.
\end{solution}

\begin{exercise}\label{exr:vii42-07}
Give an analytic surface that is umbilic everywhere.
\end{exercise}

\begin{hint}
Think constant positive curvature.
\end{hint}

\begin{solution}
A round sphere.
\end{solution}

\begin{exercise}\label{exr:vii42-08}
Give the two principal-direction families on a circular cylinder.
\end{exercise}

\begin{hint}
One bends; one does not.
\end{hint}

\begin{solution}
The circumferential circles and the axial generators.
\end{solution}

\begin{exercise}\label{exr:vii42-09}
Define \(D_{e_i}k_i\).
\end{exercise}

\begin{hint}
It is a directional derivative.
\end{hint}

\begin{solution}
\(D_{e_i}k_i=\langle\nabla k_i,e_i\rangle\).
\end{solution}

\begin{exercise}\label{exr:vii42-10}
Under the convention of this chapter, what first-order equation defines a \(k_1\)-ridge candidate?
\end{exercise}

\begin{hint}
Stationarity along the first principal direction.
\end{hint}

\begin{solution}
\(D_{e_1}k_1=0\).
\end{solution}

\begin{exercise}\label{exr:vii42-11}
What second-order sign selects a strict \(k_1\)-ridge in this convention?
\end{exercise}

\begin{hint}
A ridge is a local maximum along \(e_1\).
\end{hint}

\begin{solution}
\(D_{e_1}^2k_1<0\).
\end{solution}

\begin{exercise}\label{exr:vii42-12}
What equations select a strict \(k_2\)-valley in this convention?
\end{exercise}

\begin{hint}
Use a local minimum along \(e_2\).
\end{hint}

\begin{solution}
\(D_{e_2}k_2=0\) and \(D_{e_2}^2k_2>0\).
\end{solution}

\begin{exercise}\label{exr:vii42-13}
Why must ridge/valley code document its normal convention?
\end{exercise}

\begin{hint}
Reverse the normal.
\end{hint}

\begin{solution}
Reversing the normal flips the shape operator and principal-curvature signs, which can change signed ridge/valley naming and curvature ordering.
\end{solution}

\begin{exercise}\label{exr:vii42-14}
Are ridge curves necessarily lines of curvature?
\end{exercise}

\begin{hint}
Compare an implicit locus with an integral curve.
\end{hint}

\begin{solution}
No.  A ridge is a zero/extremum locus of a directional curvature derivative; its tangent need not equal the principal direction used in that derivative.
\end{solution}

\begin{exercise}\label{exr:vii42-15}
What local polynomial model is commonly used to estimate curvature on a mesh?
\end{exercise}

\begin{hint}
Use a tangent-frame height graph.
\end{hint}

\begin{solution}
A local Monge patch, for example \(z\approx\frac12(ax^2+2bxy+cy^2)+\cdots\).
\end{solution}

\begin{exercise}\label{exr:vii42-16}
Why is one-ring curvature estimation often noisy?
\end{exercise}

\begin{hint}
Curvature uses second derivatives.
\end{hint}

\begin{solution}
Second derivatives amplify positional and normal noise, and a small irregular neighborhood provides little averaging or conditioning.
\end{solution}

\begin{exercise}\label{exr:vii42-17}
Give a practical numerical test for a near-umbilic vertex.
\end{exercise}

\begin{hint}
Compare the two estimated principal curvatures.
\end{hint}

\begin{solution}
Mark it near-umbilic when \(|k_1-k_2|<\varepsilon_k\) for a scale-aware tolerance.
\end{solution}

\begin{exercise}\label{exr:vii42-18}
Why are principal directions ill-conditioned near an umbilic?
\end{exercise}

\begin{hint}
Eigenvalues nearly coincide.
\end{hint}

\begin{solution}
When the eigenvalue gap is small, small perturbations of the estimated shape operator can cause large rotations of its eigenvectors.
\end{solution}

\begin{exercise}\label{exr:vii42-19}
How can neighboring representatives of the same principal line be locally sign-aligned?
\end{exercise}

\begin{hint}
Use a dot product.
\end{hint}

\begin{solution}
Replace one vector by its negative when necessary so its dot product with the chosen neighboring representative is nonnegative or maximal.
\end{solution}

\begin{exercise}\label{exr:vii42-20}
How can a mesh locate first-order ridge candidates from a sampled derivative scalar \(r_1\)?
\end{exercise}

\begin{hint}
Look along edges.
\end{hint}

\begin{solution}
Detect sign changes of \(r_1\) across edges and interpolate zero crossings within incident triangles.
\end{solution}

\begin{exercise}\label{exr:vii42-21}
Why is a derivative zero crossing alone insufficient for a stable ridge?
\end{exercise}

\begin{hint}
Noise can create many sign changes.
\end{hint}

\begin{solution}
It must be filtered by the appropriate second-derivative sign, feature strength, confidence, scale persistence, and singularity handling.
\end{solution}

\begin{exercise}\label{exr:vii42-22}
Give a line-field-aware angular error between estimated and exact unit directions.
\end{exercise}

\begin{hint}
Remove the sign ambiguity.
\end{hint}

\begin{solution}
\(\theta=\arccos(|\langle e,e_*\rangle|)\).
\end{solution}

\begin{exercise}\label{exr:vii42-23}
Why does mean curvature alone not determine principal directions?
\end{exercise}

\begin{hint}
Mean curvature is only one scalar combination.
\end{hint}

\begin{solution}
It supplies essentially \((k_1+k_2)/2\) but not the eigenvectors or the separation of the two principal curvatures; a fuller shape-operator estimate is required.
\end{solution}

\begin{exercise}\label{exr:vii42-24}
What global subject naturally follows the singular line-field phenomena seen here?
\end{exercise}

\begin{hint}
The next volume studies global deformation and invariants.
\end{hint}

\begin{solution}
Algebraic topology: homotopy and homology provide tools for understanding global topological constraints beyond local differential geometry.
\end{solution}


\section{Solved problem dossiers}

\begin{problem}[Euler extrema]\label{prob:vii42-01}
Using Euler's formula with \(k_1>k_2\), show that the maximum and minimum normal curvatures occur in the principal directions.
\end{problem}

\begin{solution}
Write \(k_n(\theta)=k_2+(k_1-k_2)\cos^2\theta\).  Since \(0\le\cos^2\theta\le1\), the maximum \(k_1\) occurs for \(\theta=0\) modulo \(\pi\), and the minimum \(k_2\) for \(\theta=\pi/2\) modulo \(\pi\).
\end{solution}

\begin{problem}[Sphere regression]\label{prob:vii42-02}
A curvature estimator on a round sphere reports stable, highly anisotropic principal directions.  Why is this suspicious?
\end{problem}

\begin{solution}
A sphere is umbilic everywhere, so the principal eigenvalue gap should vanish.  The directions are mathematically undetermined; stable anisotropy is likely coming from triangulation bias, noise, or the estimator rather than the smooth surface.
\end{solution}

\begin{problem}[Cylinder ridge test]\label{prob:vii42-03}
On an ideal cylinder, \(k_1=1/R\) and \(k_2=0\) are constant.  Why does the first-order ridge equation hold everywhere, and why should a feature extractor not label the entire cylinder a sharp ridge?
\end{problem}

\begin{solution}
All directional derivatives of the constant principal curvatures vanish, so the stationarity equations hold identically.  But the second derivative also vanishes and there is no strict local extremum; strength and strictness criteria prevent the whole cylinder from becoming a ridge.
\end{solution}

\begin{problem}[Normal reversal]\label{prob:vii42-04}
Suppose an algorithm uses outward normals and finds principal curvatures \(k_1,k_2\).  What happens after all normals are reversed?
\end{problem}

\begin{solution}
With \(S=-dN\), reversing \(N\) changes \(S\) to \(-S\), so the same principal lines remain but both eigenvalues change sign.  If the code reorders eigenvalues afterward, their labels may also swap.
\end{solution}

\begin{problem}[Line-angle metric]\label{prob:vii42-05}
An exact principal direction is \(e_*=(1,0)\), while an estimator returns \(e=(-\cos5^\circ,-\sin5^\circ)\).  What is the line-field angular error?
\end{problem}

\begin{solution}
The sign is irrelevant.  The absolute dot product is \(\cos5^\circ\), so \(\theta_{\rm line}=5^\circ\), not \(175^\circ\).
\end{solution}

\begin{problem}[Near-umbilic conditioning]\label{prob:vii42-06}
Two estimated principal curvatures are \(1.0000\) and \(0.9999\).  Explain why a principal-direction feature may be unreliable even if the curvature values themselves look accurate.
\end{problem}

\begin{solution}
The eigenvalue gap is only \(10^{-4}\).  Eigenvectors of a symmetric matrix become highly sensitive to perturbations when eigenvalues nearly coincide, so the direction can rotate substantially under tiny fitting errors.
\end{solution}

\begin{problem}[Monge patch eigenvectors]\label{prob:vii42-07}
A small-slope quadratic fit has Hessian \(\begin{pmatrix}3&0\\0&1\end{pmatrix}\).  What are the approximate principal curvatures and directions in the local frame?
\end{problem}

\begin{solution}
The eigenvalues are approximately \(3\) and \(1\), with eigenvectors along the local \(x\)- and \(y\)-axes, respectively.
\end{solution}

\begin{problem}[Sign-coherent interpolation]\label{prob:vii42-08}
Two adjacent vertices store eigenvectors \(e\) and approximately \(-e\) for the same smooth principal line.  What goes wrong if they are linearly interpolated without sign alignment?
\end{problem}

\begin{solution}
They can cancel near the midpoint, creating an artificial near-zero vector and an unstable normalized direction.  Flipping one representative first avoids the spurious cancellation locally.
\end{solution}

\begin{problem}[Zero-crossing interpolation]\label{prob:vii42-09}
Along an edge, a ridge derivative estimate has values \(r(a)=-2\) and \(r(b)=6\).  Under linear interpolation, at what fraction from \(a\) does the zero occur?
\end{problem}

\begin{solution}
Solve \(-2+t(6-(-2))=0\), giving \(t=2/8=1/4\).  The candidate lies one quarter of the way from \(a\) to \(b\).
\end{solution}

\begin{problem}[Multiscale persistence]\label{prob:vii42-10}
A feature appears at one-ring scale but disappears for every moderately larger fitting radius.  How should this affect confidence?
\end{problem}

\begin{solution}
Confidence should decrease.  A feature that exists only at the smallest, noise-sensitive scale is less likely to represent persistent surface geometry than one that survives across a reasonable scale interval.
\end{solution}

\begin{problem}[Laplacian versus principal data]\label{prob:vii42-11}
A program knows only a reliable mean-curvature normal vector at each vertex.  Can it reconstruct both principal directions uniquely?
\end{problem}

\begin{solution}
Not in general.  Mean curvature gives only the sum of principal curvatures (with normal direction) and lacks the anisotropic second-fundamental-form information needed to recover the two eigenlines.
\end{solution}

\begin{problem}[Tracing stop condition]\label{prob:vii42-12}
A curvature-line tracer approaches a region where \(|k_1-k_2|\) falls below tolerance.  Why is stopping or reducing confidence preferable to forcing a direction through it?
\end{problem}

\begin{solution}
The principal line becomes ill-conditioned near an umbilic.  Forcing a direction lets numerical noise choose an arbitrary continuation and can create visually plausible but geometrically unsupported curves.
\end{solution}


\section{Challenge problems}

\begin{problem}[Orthogonality of curvature lines]\label{prob:vii42-ch01}
Prove that at a nonumbilic point the two curvature-line families meet orthogonally, starting only from self-adjointness of the shape operator.
\end{problem}

\begin{solution}
Let \(Se_1=k_1e_1\) and \(Se_2=k_2e_2\) with \(k_1\ne k_2\).  Self-adjointness gives \(k_1\langle e_1,e_2\rangle=\langle Se_1,e_2\rangle=\langle e_1,Se_2\rangle=k_2\langle e_1,e_2\rangle\).  Hence \((k_1-k_2)\langle e_1,e_2\rangle=0\), so the inner product vanishes.
\end{solution}

\begin{problem}[Why ridge tangents differ]\label{prob:vii42-ch02}
Let \(r=D_{e_1}k_1\).  At a regular ridge point where \(r=0\) and \(\nabla r\ne0\), what determines the tangent to the ridge locus, and why need it not be \(e_1\)?
\end{problem}

\begin{solution}
The ridge is locally the level set \(r=0\), so its tangent vectors are orthogonal to \(\nabla r\).  The defining principal direction \(e_1\) appears inside the scalar \(r\), but there is no general reason for \(e_1\) to be orthogonal to \(\nabla r\); hence the ridge tangent need not equal \(e_1\).
\end{solution}

\begin{problem}[Perturbation near an umbilic]\label{prob:vii42-ch03}
Consider symmetric matrices \(S_\varepsilon=\begin{pmatrix}1&\varepsilon\\\varepsilon&1\end{pmatrix}\).  Compare their eigenvectors with those of \(S_0=I\) and explain the numerical lesson for principal-direction tracing.
\end{problem}

\begin{solution}
For every nonzero \(\varepsilon\), the eigenvectors are the diagonal directions \((1,1)\) and \((1,-1)\), while for \(\varepsilon=0\) every direction is an eigenvector.  An arbitrarily small perturbation selects specific directions out of a completely degenerate eigenspace, illustrating why directions near umbilics are intrinsically unstable.
\end{solution}

\begin{problem}[Orientation-invariant feature proposal]\label{prob:vii42-ch04}
Design an orientation-insensitive scalar candidate for strong anisotropic bending using the principal curvatures, and explain one advantage and one limitation compared with signed ridge/valley labels.
\end{problem}

\begin{solution}
One example is \(|k_1-k_2|\), or a normalized variant such as \(|k_1-k_2|/(|k_1|+|k_2|+\varepsilon)\).  It is unchanged when the normal is reversed and highlights anisotropy, but it does not distinguish maxima from minima along principal directions and therefore is not itself a signed ridge/valley definition.
\end{solution}

\begin{problem}[End-to-end validation design]\label{prob:vii42-ch05}
Design an end-to-end regression suite for curvature-line and ridge extraction that separates estimator error from tracer error and from feature-selection error.
\end{problem}

\begin{solution}
First use analytic surfaces to compare estimated \(k_i\) and principal lines against ground truth using sign-invariant angular error.  Next feed exact or synthetic principal line fields to the tracer to isolate integration and face-crossing behavior.  Then feed analytic curvature and derivative samples to zero-crossing/selection code to test ridge logic.  Finally run the complete pipeline over meshes at multiple resolutions, noise levels, and scales, checking sphere umbilics, cylinder line families, known torus/Monge features, continuity, and persistence.
\end{solution}


\section{Volume VII synthesis}

Volume VII began with topological and smooth manifolds, built tangent and
cotangent geometry, developed bundles and differential forms, passed through the
classical geometry of curves and surfaces, and then constructed Riemannian,
Lorentzian, and hyperbolic geometry.  The final computational arc translated
geodesic distance, heat flow, Laplace--Beltrami operators, and curvature features
into algorithms on meshes.

The recurring theme is that local differential data and global structure must be
kept distinct but compatible.  Metrics produce lengths and geodesics; connections
produce parallel transport and curvature; shape operators produce principal
geometry; discrete operators approximate these objects only after choices of
mesh, mass, scale, and numerical convention are made explicit.

\section{Bridge to Volume VIII}

Principal line fields become singular at umbilics, geodesics interact with
nontrivial topology, and global geometry cannot be understood from local
coordinates alone.  Volume VIII begins with homotopies of maps and develops the
algebraic-topological invariants that measure global deformation, loops, cells,
and homology.  The passage is therefore natural: differential geometry describes
local geometric structure, while algebraic topology organizes the global
constraints that local geometry must obey.
